\documentclass{ximera}

\title{Exponents and Positive Powers of 10: Summary}   % Use xmlatex all -s to compile when SERVE refuses to work.
\author{Kelly Stady}
\license{CC: 0}         % replace with an appropriate license, or set it in xmPreamble

\begin{document}

\begin{abstract}

\end{abstract}

\maketitle
\label{xim:ExponentsPowers10Summary}


\section*{Summary}

\normalsize
\begin{enumerate}
\item \textbf{Exponential Notation}

Exponential notation is shorthand for repeated multiplication of the same quantity.  A number written in exponential notation has a base and an exponent.  The exponent indicates the number of times to multiply the base.  For example, the product $8 \cdot 8 \cdot 8$ is written as $8^3$ in exponential notation.  The number 8 is the base and 
the number 3 is the exponent. 

\item \textbf{Positive Powers of 10}

A \textbf{positive power of 10} is a power of 10 in which the exponent is a positive integer: $\{1, 2, 3, 4, \ldots \}.$  Below are the first four positive powers of 10.
\begin{eqnarray*}
10^1 &=& 10 \\
10^2 &=& 100   \\
10^3 &=& 1,000 \\
10^4 &=& 10,000
\end{eqnarray*}

In general, if $n$ is a positive integer, then the value of $10^n$ is the number 1 followed by $n$ zeros.

\item \textbf{Multiplying a Number by a Positive Power of 10}

To multiply a number by a positive power of 10, count the number of zeros in the power of 10, then move the decimal \\ point to the \textbf{right} that many places.


\item \textbf{Dividing a Number by a Positive Power of 10}

To divide a number by a positive power of 10, count the number of zeros in the power of 10, then move the decimal \\ point to the \textbf{left} that many places.


\end{enumerate}


\end{document}
